\appendix

\section{Formal Development of the Security Preservation}\label{app:preservation}

The body (Section~\ref{sec:security}) states the security-preservation theorem
(Theorem~\ref{thm:security-preserve}) and sketches its proof. This appendix gives
the formal development for audit: the building-block properties the reduction
treats as black boxes, each carrying an error term; the relation-preservation
lemma; the step-by-step reductions; a corollary for the case in which the open
premises are closed; and a remark relating the result to BDEC. Throughout,
$\mathcal A$ denotes a classical probabilistic polynomial-time adversary and all
probabilities are in the random-oracle model. The advantages
$\mathsf{Adv}^{\bullet}_{\mathsf{zkVM\text{-}BDEC}}$ are those of
Definition~\ref{def:zkvm-bdec-adv}.

\subsection{Building-block properties}\label{app:prsv-defs}

Each property is stated as the existence of an algorithm together with an error
term. The premises of Theorem~\ref{thm:security-preserve} assert the relevant
errors negligible; we mark which properties are assumed and which carry an open
premise.

\begin{definition}[Knowledge soundness]\label{def:ks}
The zkVM argument $\Pi$ is \emph{knowledge-sound with knowledge error
$\epsilon_{\mathrm{KS}}$} if there is a probabilistic polynomial-time extractor
$\mathcal{E}$ with oracle access to the prover such that, for every prover
$P^\ast$ that makes $\Pi.\mathsf{Verify}(x,\cdot)$ accept with probability
$\varepsilon$, the extractor $\mathcal{E}^{P^\ast}(x)$ outputs a witness $w$ with
$R(x;w)=1$ with probability at least $\varepsilon-\epsilon_{\mathrm{KS}}$. This is
the standard knowledge-soundness-with-error formulation; $\epsilon_{\mathrm{KS}}$
is the additive extraction-failure term that enters the unforgeability
bound~\eqref{eq:adv-unf}, and the concrete soundness error of the STARK and BCS
arguments we rely on is treated in \S\ref{app:argument}. (Assumed; classical
random-oracle model.)
\end{definition}

\begin{definition}[Arithmetisation soundness error]\label{def:eair}
$$\epsilon_{\mathrm{AIR}}=\Pr[\text{the Griffin AIR accepts a trace whose committed
transition differs from }\pi]$$, 
the probability taken over the prover-chosen trace
and the argument verifier's randomness, where $\pi$ is the reference permutation of
Definition~\ref{def:ref-perm}. Premise~(T2) asserts $\epsilon_{\mathrm{AIR}}$
negligible; the supporting argument is Appendix~\ref{app:proofs}.
\end{definition}

\begin{definition}[Signature-transcript simulatability]\label{def:szk}
PLUM \emph{admits a signature-transcript simulator with error
$\epsilon_{\mathrm{sZK}}$} if there is a probabilistic polynomial-time algorithm
$\mathcal{S}_\sigma$ such that, for every message $m$ and public key, the output of
$\mathcal{S}_\sigma$ run without the secret key is
$\epsilon_{\mathrm{sZK}}$-indistinguishable from a credential or pseudonym
transcript produced with the secret key. Premise~(T4) asserts
$\epsilon_{\mathrm{sZK}}$ negligible. This premise is \emph{open}: PLUM's analysis
establishes EU-CMA security, which does not by itself furnish such a simulator.
\end{definition}

\begin{definition}[Zero knowledge]\label{def:zk}
The zkVM argument $\Pi$ is \emph{zero-knowledge with error
$\epsilon_{\mathrm{ZK}}$} if there is a probabilistic polynomial-time simulator
$\mathcal{S}_\Pi$ such that, for every $(x,w)$ with $R(x;w)=1$, the simulated proof
$\mathcal{S}_\Pi(x)$ is $\epsilon_{\mathrm{ZK}}$-indistinguishable from an honest
proof $\Pi.\mathsf{Prove}(x,w)$. (Assumed for the analysis. In the measured deployment this premise is
not met: the produced receipt mode (the default \texttt{prove} path) has no
established ZK simulator, and the ZK-wrapped mode
(\texttt{.compressed().groth16()}) is not produced within the resource and
memory bound; see Section~\ref{ssec:zk-gap}.)
\end{definition}

\begin{definition}[Existential unforgeability]\label{def:euf}
PLUM is \emph{EU-CMA secure with error $\epsilon_{\mathrm{EUF}}$} if every
probabilistic polynomial-time forger, given the public key and a signing oracle,
produces a valid signature on a message it never queried with probability at most
$\epsilon_{\mathrm{EUF}}$. (Established in PLUM's analysis,
\cite[Theorem~1]{10.1007/978-981-95-2961-2_6}, in the classical random-oracle
model.)
\end{definition}

\subsection{Relation preservation}\label{app:prsv-relation}

\begin{proof}[Proof of Lemma~\ref{lem:relation-preserve}]
We account for the operations a guest evaluation invokes. The credential-generation
relation $R^{\Sig}_{\mathsf{cre}}$, two PLUM verifications under $\pk_U$, decomposes
into large-field multiplications, returned identically by the
modular-multiplication precompile, and Griffin permutations inside PLUM's hashing
and inside every Merkle-path and low-degree-test recomputation, returned
identically by the Griffin precompile under the lookup binding, together with
native control flow, whose faithful execution follows from the \emph{soundness} of
the zkVM argument for the RISC-V execution relation (an accepting proof certifies a
valid execution); the stronger \emph{knowledge}-soundness (Definition~\ref{def:ks})
that extracts a witness is invoked separately, by the unforgeability reduction. For
$R^{\Sig}_{\mathsf{cre}}$ this enumeration is complete at the predicate level, conditional on the open statement-binding refinement of Section~\ref{sec:system} and on premise~(T2). The presentation relation
$R^{\Sig}_{\mathsf{show}}$ adds only $k+2$ further signature-verification
predicates built from the same multiplications and Griffin permutations; its
ledger-membership, non-revocation, and disclosure checks are performed by the
relying party outside the arithmetised relation. Composing the per-operation
equalities along the finite, input-fixed execution yields equality of the
accept-or-reject decision at the program level. This program-level equality lifts
to credential-verifier acceptance only when the zkVM journal commits the public
statement $x$; the specification assumes this binding, whereas the current
RISC~Zero prototype commits the verification outcome but not $x_{\mathsf{cre}}$
itself (Section~\ref{sec:system}), so the code-level reading is conditional on
that journal fix. The Griffin step's constraint-level
soundness rests on Appendix~\ref{app:proofs}, which is argued rather than
machine-verified, so the equality holds modulo premise~(T2).
\end{proof}

\subsection{Proof of the preservation theorem}\label{app:prsv-thm}

We re-instantiate BDEC's three reductions~\cite{10.1007/978-981-96-0957-4_3} with
PLUM as the signature scheme $\Sig$ and the zkVM as the zkSNARK $\Pi$; by
Lemma~\ref{lem:relation-preserve} the two provers accept the same relation.

\begin{proof}[Proof of Theorem~\ref{thm:security-preserve}, unforgeability]
Let $\mathcal A$ win the zkVM-BDEC unforgeability experiment. We build a PLUM
EU-CMA forger $\mathcal B$. $\mathcal B$ runs $\mathcal A$, answering its issuance
and showing queries with its own signing oracle and the zkVM prover, so
$\mathcal A$'s view matches the real experiment. When $\mathcal A$ outputs an
accepting $(x^\star,\pi^\star)$ for a statement outside the issued set, $\mathcal B$
invokes the extractor $\mathcal{E}$ of Definition~\ref{def:ks} to obtain
$w^\star=(\pk_U^\star,\mathsf{psk}^\star)$, which fails with probability at most
$\epsilon_{\mathrm{KS}}$. Conditioned on extraction, $w^\star$ satisfies
$R^{\Sig}_{\mathsf{cre}}$ except when the committed Griffin transition differs from
the reference permutation, an event of probability at most $\epsilon_{\mathrm{AIR}}$
by Definition~\ref{def:eair}; otherwise the extracted $\mathsf{psk}^\star$ (the pseudonym secret, structurally a PLUM signature under $\pk_U$, Section~\ref{sec:prelim}) is valid
on a message $\mathcal A$ never had signed, and by premise~(T3) the
precompile substitution leaves PLUM's EU-CMA reduction unchanged, so
$\mathsf{psk}^\star$ is a forgery against the unmodified scheme, which $\mathcal B$
outputs.
Writing $\epsilon_{\mathrm{EUF}}$ for PLUM's EU-CMA advantage bound,
\[
  \mathsf{Adv}^{\mathsf{unf}}_{\mathsf{zkVM\text{-}BDEC}}(\mathcal A)
  \le \epsilon_{\mathrm{EUF}}+\epsilon_{\mathrm{KS}}+\epsilon_{\mathrm{AIR}},
\]
which is~\eqref{eq:adv-unf}.
\end{proof}

\begin{proof}[Proof of Theorem~\ref{thm:security-preserve}, anonymity]
Let $\mathcal A$ be the anonymity distinguisher, with advantage
$\bigl|\Pr[b'=b]-\tfrac12\bigr|$. We proceed through a sequence of games on the
challenge showing for the credential $\mathsf{cred}_b$. Game~$\mathsf{G}_0$ is the
real experiment. In game~$\mathsf{G}_1$ the pseudonym key and credential are
produced by $\mathcal{S}_\sigma$ rather than with the secret key, so by
Definition~\ref{def:szk}
$\bigl|\Pr[\mathsf{G}_0{=}1]-\Pr[\mathsf{G}_1{=}1]\bigr|\le\epsilon_{\mathrm{sZK}}$.
The role of $\mathcal{S}_\sigma$ is to let the reduction produce the challenge
credential $c_b$ and pseudonym signature $\mathsf{psk}_b$ without the challenge
user's long-term secret key $\sk_U$, matching BDEC's anonymity simulator
(\cite[Appendix~A.2]{10.1007/978-981-96-0957-4_3}); the hop is load-bearing
whenever PLUM's transcripts cannot be produced without $\sk_U$.
In game~$\mathsf{G}_2$ the proof is produced by $\mathcal{S}_\Pi$ rather than from
the witness, so by Definition~\ref{def:zk}
$\bigl|\Pr[\mathsf{G}_1{=}1]-\Pr[\mathsf{G}_2{=}1]\bigr|\le\epsilon_{\mathrm{ZK}}$.
In $\mathsf{G}_2$ we account for every $b$-dependent component of the challenge
statement
$x_{\mathsf{show}}=(A^\downarrow,(\mathsf{ppk}_{U,TA}^{(j)})_{j=1}^{k},\mathsf{ppk}_{U,V})$.
The disclosed attribute set $A^\downarrow$ is equalised across the two credentials
by the well-formedness condition
$\mathsf{Leak}(\mathsf{stmt}_0)=\mathsf{Leak}(\mathsf{stmt}_1)$. The verifier-facing
pseudonym $\mathsf{ppk}_{U,V}$ is sampled freshly at showing time,
$\mathsf{ppk}_{U,V}\leftarrow\{0,1\}^{\lambda}$, independently of which credential
is shown, so its distribution is uniform regardless of $b$. The teaching-authority
pseudonyms $(\mathsf{ppk}_{U,TA}^{(j)})_{j}$ are issuance-time, credential-bound
values that the well-formedness condition equalises across the two challenge
credentials, so $\mathsf{Leak}$ includes the multiset $\{\mathsf{ppk}_{U,TA}^{(j)}\}$
(Section~\ref{sec:prelim}). The secret keys, the long-term key $\pk_U$, and the
credential $c_{U,V}$ are witness-only: they were replaced by $\mathcal{S}_\sigma$
in $\mathsf{G}_1$ and removed from the proof by $\mathcal{S}_\Pi$ in $\mathsf{G}_2$.
Every $b$-dependent component being equalised, freshly sampled, or simulated, the
challenge is distributed identically under $b=0$ and $b=1$, so
$\Pr[b'=b]=\tfrac12$.
Summing the two hops,
\[
  \mathsf{Adv}^{\mathsf{anon}}_{\mathsf{zkVM\text{-}BDEC}}(\mathcal A)
  \le \epsilon_{\mathrm{ZK}}+\epsilon_{\mathrm{sZK}},
\]
which is~\eqref{eq:adv-anon}. No witness is extracted, so $\epsilon_{\mathrm{KS}}$
does not enter; the challenge showing is honestly generated, so
$\epsilon_{\mathrm{AIR}}$ does not enter.
\end{proof}

\begin{proof}[Proof of Theorem~\ref{thm:security-preserve}, unlinkability]
BDEC's unlinkability experiment
(\cite[Appendix~A.3]{10.1007/978-981-96-0957-4_3}) forwards a single challenge
credential $c^{(b)}$ issued under one of two pseudonym keys $\mathsf{ppk}^{(b)}$,
with the verifier-facing pseudonym freshly sampled as in the anonymity step.
BDEC's A.3 simulator simulates the pseudonym signature and the credential with
$\mathcal{S}_\sigma$ and the proof with $\mathcal{S}_\Pi$. Applying the same two
hops to the single transcript, $\mathcal{S}_\sigma$ for
$(\mathsf{psk}^{(b)},c^{(b)})$ (cost $\epsilon_{\mathrm{sZK}}$) and
$\mathcal{S}_\Pi$ for the proof (cost $\epsilon_{\mathrm{ZK}}$), makes it
independent of $b$, so
\[
  \mathsf{Adv}^{\mathsf{unlink}}_{\mathsf{zkVM\text{-}BDEC}}(\mathcal A)
  \le \epsilon_{\mathrm{ZK}}+\epsilon_{\mathrm{sZK}}.
\]
The single-transcript game incurs each simulation once; there is no factor of two.
One difference of presentation: BDEC's A.3 carries $c^{(b)}$ as a public statement
component, so its $\mathcal{S}_\sigma$ hop is forced; in this thesis's relation
$R^{\Sig}_{\mathsf{show}}$ the credential $c_{U,V}$ is a witness component
(Section~\ref{sec:prelim}), so the $\mathcal{S}_\Pi$ hop already hides it and
$\mathcal{S}_\sigma$ is needed only to produce $c^{(b)}$ without $\sk_U$ inside the
reduction. Either way $\epsilon_{\mathrm{sZK}}$ enters once.
\end{proof}

\begin{corollary}[Preservation under closed premises]\label{cor:unconditional}
Suppose premises~(T2) and~(T4) hold, so that $\epsilon_{\mathrm{AIR}}$ and
$\epsilon_{\mathrm{sZK}}$ are negligible; the zkVM is knowledge-sound and
zero-knowledge with negligible $\epsilon_{\mathrm{KS}},\epsilon_{\mathrm{ZK}}$; and
PLUM is EU-CMA secure. Then zkVM-BDEC is unforgeable, anonymous, and unlinkable.
Against quantum adversaries the same conclusion holds once the
quantum-random-oracle lifts of PLUM's EU-CMA reduction and of the nested
Fiat--Shamir compilation are established.
\end{corollary}

\begin{proof}
Each bound of Theorem~\ref{thm:security-preserve} is a sum of the named error
terms; under the stated hypotheses every term is negligible, so each advantage is
negligible. The quantum clause adds the two open lifts as further hypotheses.
\end{proof}

\begin{remark}[Relation to BDEC]\label{rem:reinstantiation}
Theorem~\ref{thm:security-preserve} re-instantiates BDEC's Theorems~1--3, proved
for a generic $(\text{signature},\text{zkSNARK})$ pair, with the concrete pair
$(\text{PLUM},\text{zkVM})$. The contribution is not a new reduction but the
explicit accounting of the substrate-specific errors $\epsilon_{\mathrm{AIR}}$ (the
precompile's arithmetisation soundness) and
$\epsilon_{\mathrm{KS}},\epsilon_{\mathrm{ZK}}$ (the zkVM's argument), together with
the identification of $\epsilon_{\mathrm{sZK}}$ as the open premise~(T4) that BDEC's
generic anonymity argument silently requires of the signature. We further note
that the relation $R^{\Sig}_{\mathsf{show}}$ places the credential $c_{U,V}$ in the
witness rather than, as in BDEC's literal layout, the public statement; the result
is therefore a re-derivation of BDEC's Theorems~1--3 for this concrete pair, not a
verbatim re-instantiation.
\end{remark}

\section{Evidence for Premise~(T2): Specification-Level Soundness of the Griffin-$\mathbb{F}_p^{192}$ AIR}\label{app:proofs}

This appendix develops the soundness argument for the Griffin-$\mathbb{F}_p^{192}$
permutation AIR of Section~\ref{sec:system}, on which the conditional
precompile-soundness theorem (Theorem~\ref{thm:precompile-sound}) and the
relation-preservation lemma of Section~\ref{sec:security} rest. We follow the
standard practice for arithmetisation correctness. As in
ethSTARK~\cite{ethstark} (Definitions~3 and~4), the verified Cairo
representation~\cite{cairoverify}, and per-gadget AIR-soundness arguments such
as~\cite{divremair}, the formal object is a constraint-soundness statement, a
universally quantified implication over satisfying traces, together with a
separate completeness statement, while the knowledge-soundness of the surrounding
argument is treated apart in Section~\ref{app:argument}. The malicious prover is
captured by the quantifier ``for every trace that satisfies the constraints'',
so the statement uses no challenger, oracle, or advantage function. The one component that
is supported empirically rather than proved, namely that the implemented chip
realises the constraint set analysed here, appears as the explicit
premise of the soundness theorem and is discussed in Section~\ref{app:status}.

\subsection{Setup}\label{app:air-setup}

\begin{definition}[Reference permutation]\label{def:ref-perm}
Let $\pi\colon\mathbb{F}_p^{4}\to\mathbb{F}_p^{4}$ denote one application of the
Griffin permutation at PLUM's parameters, a degree-$3$ S-box over a width-$4$
state with $14$ rounds and the SHAKE256-derived round constants and base quadratic
coefficients $(\alpha,\beta)$, scaled per lane by Griffin's index law, fixed in
Section~\ref{sec:prelim}, as computed by the reference implementation.
\end{definition}

\begin{definition}[Algebraic satisfiability]\label{def:air-sat}
The Griffin AIR is a finite constraint set $\mathcal{C}=\{(Q_i,H_i)\}_i$ over a
trace $w$ comprising $14$ per-round rows and one controller row, where each $Q_i$
is a constraint polynomial and $H_i$ its enforcement domain. A trace $w$
\emph{satisfies} $\mathcal{C}$ on input $s\in\mathbb{F}_p^{4}$ with claimed output
$s'\in\mathbb{F}_p^{4}$ if its first-row input lanes decode to $s$, its last-row
output lanes decode to $s'$, and every constraint polynomial $Q_i$ vanishes on its
enforcement domain $H_i$ under $w$.
\end{definition}

\subsection{Soundness}\label{app:air-soundness}

\begin{theorem}[Specification-level AIR soundness]\label{thm:air-sound}
Suppose the implemented chip, comprising the per-round chip together with its
sibling controller chip, enforces the constraint set $\mathcal{C}$ analysed below. Then for every input $s\in\mathbb{F}_p^{4}$ and every trace $w$
satisfying $\mathcal{C}$ with claimed output $s'$, we have $s'=\pi(s)$.
Equivalently, $\mathcal{C}$ admits no satisfying trace whose committed output
differs from the reference permutation.
\end{theorem}

\begin{proof}
By Lemmas~\ref{lem:round-fn} and~\ref{lem:compose}. Lemma~\ref{lem:round-fn}
fixes, for each round, the non-linear, linear, and round-constant layers of $w$ to
equal the reference round's outputs on its inputs, and Lemma~\ref{lem:compose}
forces the $14$ rows to compose as $14$ sequential rounds on a single state, bound
on the controller row to the syscall's memory input and output. Composing the
per-round equalities of Lemma~\ref{lem:round-fn} along the schedule fixed by
Lemma~\ref{lem:compose} yields $s'=\pi(s)$. The composition is unconditional given
the two lemmas, and the status of the chip--constraint equality (premise~(T2)) is discussed
in Section~\ref{app:status}.
\end{proof}

\subsection{Per-family lemmas}\label{app:air-families}

The ten constraint families group into the round function (B-1 through B-5) and
the schedule, binding, and canonicality apparatus (B-6 through B-10). Each clause
states the reference relation a family enforces and, in contrapositive, the
constraint a deviating trace must violate.

\begin{lemma}[Round-function soundness; families B-1 through B-5]\label{lem:round-fn}
Fix a round with input lanes $s=(s_0,s_1,s_2,s_3)$ and let $s'$ be the trace's
claimed round output. Under constraint families B-1 through B-5, every trace
satisfying $\mathcal{C}$ has $s'=\rho(s)$, where $\rho$ denotes one reference
Griffin round.
\end{lemma}

\begin{proof}
For each family we exhibit the constraint polynomial it enforces on the round
domain and observe that any deviating value leaves that polynomial non-vanishing,
hence unsatisfying; equality of the surviving assignment with $\rho(s)$ then
follows by composing the layers.
\emph{Forward S-box} (lane $1$, family B-1). The chip enforces
$Q_{\mathsf{sq}}=\mathit{sq}-s_1^{2}$ and
$Q_{\mathsf{cube}}=\mathit{cube}-\mathit{sq}\cdot s_1$, both vanishing on the round
domain, so the lane-$1$ nonlinear output is $\mathit{cube}=s_1^{3}\bmod p$; any
output $\neq s_1^{3}$ leaves $Q_{\mathsf{cube}}$ non-zero.
\emph{Inverse S-box} (lane $0$, family B-2). The chip enforces
$Q_{\mathsf{inv}}=s_0-(s^{\mathrm{nl}}_0)^{3}$, witnessing the cube root in the
backward direction; since $\gcd(3,p-1)=1$ (Remark~\ref{rem:prime-substitution},
verified for the implemented modulus in Section~\ref{sec:system}), the cube map is
a bijection on $\mathbb{F}_p$, so the witnessed $s^{\mathrm{nl}}_0$ is the unique
cube root and equals the forward inverse S-box on $s_0$.
\emph{Quadratic layer} (lanes $2,3$, family B-3). For $i\in\{2,3\}$ the chip
enforces $Q_i=s^{\mathrm{nl}}_i-s_i\,(L_i^{2}+\alpha_{i-2}L_i+\beta_{i-2})$, where
$L_i=(i-1)\,s^{\mathrm{nl}}_0+s^{\mathrm{nl}}_1+s^{\mathrm{nl}}_{i-1}$ reads the two
updated S-box lanes ($s^{\mathrm{nl}}_0$ the inverse-S-box lane,
$s^{\mathrm{nl}}_1$ the forward-S-box lane) and the preceding updated quadratic
lane $s^{\mathrm{nl}}_{i-1}$; for $i=2$ the third summand is absent, so
$L_2=s^{\mathrm{nl}}_0+s^{\mathrm{nl}}_1$. The quadratic coefficients follow
Griffin's index law: a single base pair $(\alpha,\beta)$ is read from the SHAKE256
seed, and $\alpha_{i-2}=(i-1)\alpha$, $\beta_{i-2}=(i-1)^{2}\beta$ (so
$\alpha_0=\alpha,\beta_0=\beta$ for $i=2$ and $\alpha_1=2\alpha,\beta_1=4\beta$ for
$i=3$). The base $(\alpha,\beta)$ are the first two field elements read from
$\mathrm{SHAKE256}$ seeded by the domain string $\mathrm{PlumGriffin}(p,4,2,128)$,
the same seed that produces the $52$ round constants; a wrong $\alpha,\beta$, a
wrong index scaling, or an off-by-one in $L_i$ leaves $Q_i$ non-zero.
\emph{MDS layer} (family B-4). Per output lane $j$ the chip enforces
$Q^{\mathsf{mds}}_j=s'_j-\bigl\langle \mathrm{circ}(2,1,1,1)_j,\,s^{\mathrm{nl}}\bigr\rangle$,
where $\mathrm{circ}(2,1,1,1)_j$ is the symmetric circulant row
$\mathrm{first}[(j+\mathrm{col})\bmod 4]$ matching \texttt{build\_matrix}, not the
standard $(\mathrm{col}-j)\bmod 4$ form (the two differ on rows $1,3$); any other
matrix row leaves some $Q^{\mathsf{mds}}_j$ non-zero.
\emph{Round constants} (family B-5). Per lane $\ell$ the chip enforces
$Q^{\mathsf{rc}}_\ell=s'_\ell-(\text{MDS output})_\ell-\mathrm{RC}[4r+\ell]$ under
one-hot round selection $r$, with $\mathrm{RC}$ preprocessing-bound (family B-10);
a substituted or omitted constant leaves $Q^{\mathsf{rc}}_\ell$ non-zero. The
preprocessing table holds $\mathrm{STATE\_WIDTH}\times(\text{rounds}-1)=52$
constants, applied to the first $13$ rounds; the final round applies no
round-constant layer, so its constraint is
$Q^{\mathsf{rc}}_\ell=s'_\ell-(\text{MDS output})_\ell$.

Composing B-1 through B-5 along the round, the unique satisfying assignment of
$s'$ is the reference round $\rho(s)$.
\end{proof}

\begin{lemma}[Composition, binding, canonicality; families B-6 through B-10]\label{lem:compose}
Under constraint families B-6 through B-10, every trace satisfying $\mathcal{C}$
realises its $14$ rows as the $14$ sequential reference rounds on a single state
bound to the syscall's memory input and output, with canonical output lanes.
\end{lemma}

\begin{proof}
\emph{Round count} (family B-6). The schedule is fixed to $14$ rounds by the
trace-generation row count and the controller's lookup anchors at directions $0$
and $\mathrm{NB\_ROUNDS}$, rather than by a dedicated algebraic constraint; under
the lookup binding (premise~(T1)) a trace with a different round count produces an
unbalanced lookup and has no satisfying assignment.
\emph{Memory binding} (family B-7). The controller chip constrains the input and
output state lanes to the values at the syscall's memory addresses and clock, a
memory-consistency lookup constraint resolved under premise~(T1), so a trace
reading or writing a different address or clock produces an unbalanced lookup.
\emph{Output canonicality} (family B-8). Each output lane $s'_j$ carries the range
constraint $0\le s'_j<p$ enforced by a byte-wise less-than comparison
(\texttt{FieldLtCols}); a non-canonical lane violates it.
\emph{Row scheduling} (family B-9). The real/padding and first/last-round selector
flags enforce one ordered run with cross-row state threading
$\mathsf{state\_out}_r=\mathsf{state\_in}_{r+1}$, realised by the per-row lookup
interactions (a \textsc{receive} at direction $r$ and a \textsc{send} at direction
$r{+}1$) under the lookup binding (premise~(T1)); a ghost or mis-ordered row
creates an unbalanced lookup.
\emph{Parameter immutability} (family B-10). The modulus, round count, MDS entries,
round constants, and quadratic coefficients are preprocessing-bound rather than
prover-supplied, so a per-row choice of any of them has no satisfying assignment.

By B-9, and given the lookup binding~(T1), the per-round equalities of
Lemma~\ref{lem:round-fn} chain into $14$ sequential rounds on one state; with the
round count fixed by the lookup schedule (B-6, premise~(T1)) and with B-8 and
B-10, this is the reference permutation, and with B-7 it is bound to the
syscall's memory input and output.
\end{proof}

\subsection{Completeness}\label{app:air-completeness}

\begin{lemma}[Completeness]\label{lem:air-complete}
For every input $s\in\mathbb{F}_p^{4}$, the honest Griffin execution on $s$ yields
a trace satisfying $\mathcal{C}$ with output $\pi(s)$.
\end{lemma}

\begin{proof}
Following the practice of~\cite{cairoverify}, where soundness is the
load-bearing direction and completeness is discharged operationally, we support
the claim by the cross-codebase equivalence tests: the trace
generator populated from the reference computation produces a satisfying
assignment on every tested input, namely the $256$-vector randomised suite, the
hand-picked and zero/high-entropy vectors, and the single-permutation smoke proof
that completes execute, prove, and verify (Section~\ref{sec:system}).
\end{proof}

\subsection{The argument layer}\label{app:argument}

Theorem~\ref{thm:air-sound} concerns the \emph{arithmetisation}. The surrounding
claim, that an accepting receipt implies a satisfying trace and that a witness is
extractable, is the \emph{knowledge-soundness} of the zkVM's STARK, a separate
object we do not re-establish. We take it in the IOP and extractor sense of
ethSTARK~\cite{ethstark} (Definition~4): for every prover producing an accepting
proof with non-negligible probability there exists an extractor outputting a
satisfying trace, up to a negligible knowledge error. This is the
knowledge-soundness hypothesis of the preservation theorem of
Section~\ref{sec:security}, itself obtained through the Fiat--Shamir and BCS
compilation of an interactive oracle proof. The arithmetisation soundness of this
appendix and the argument knowledge-soundness are independent objects. End-to-end
soundness needs both, and we keep them separate.

\subsection{The open step}\label{app:status}

The compositional argument of Theorem~\ref{thm:air-sound} is unconditional given
Lemmas~\ref{lem:round-fn} and~\ref{lem:compose} and the premise that the
implemented chip enforces exactly the analysed constraint set $\mathcal{C}$. That
premise is the one step we do not machine-verify. The evidence for it is the
cross-codebase agreement of the chip's native compute with the reference
permutation on $256$ random vectors, with per-vector false-pass at most
$1/p\approx2^{-199}$ and aggregate at most $2^{-191}$ across the suite, together
with the hand-picked and zero/high-entropy vectors. This per-vector bound assumes
a wrong executor produces outputs uniform over $\mathbb{F}_p$ and independent of
the reference; a deterministic structural bug agreeing on every sampled input
would not be caught, so $2^{-191}$ bounds only random disagreement on the sampled
set and is \emph{not} a bound on the open constraint-level premise. The suite pins the
\emph{executor} against the reference, but the \emph{constraint system} checked against that
executor is what the vector suite does not reach.
A separate fault-injection suite reaches the constraint system in the
\emph{rejection} direction for eight of the ten families (B-1 through B-5, B-8, and
the per-row B-9 selectors): a single-cell corruption of a satisfying trace is
rejected by the per-row constraints in each case, including the MDS lanes $1$ and
$3$ that distinguish the implemented symmetric circulant from the standard one.
This evidence is itself bounded, and we state the bounds. It is evaluated through
\texttt{debug\_constraints}, which checks the per-row algebraic constraints but not
the cross-table lookups, so the lookup-borne families (round count B-6, memory
binding B-7, and the cross-row threading of B-9) lie outside its reach and would
require a full-prover harness; parameter immutability (B-10) is enforced by
construction, with no prover-chosen column to reject. The suite thus shows that
specific local deviations are caught, not that no wrong witness satisfies
$\mathcal{C}$, which remains the open premise.
Closing it fully requires deriving the AIR witness directly from the reference code and
checking, ideally by machine, that $\mathcal{C}$ admits only that witness. We
therefore report Theorem~\ref{thm:air-sound} as argued at the specification level
and supported empirically at the implementation level, with the
chip--constraint equality carried as the open premise~(T2) of both this theorem
and the conditional precompile-soundness theorem
Theorem~\ref{thm:precompile-sound}.
